\oldpage{422}of this character, but which was not entirely satisfactory in its
action.

Dr.\ Hamilton thinks that the method of treating fractures of the clavicle
recommended by Hippocrates, and adopted by Celsus and Dupuytren,
and approved by Drs.~Eastman, Eve, Buck, and Post, of placing the
patient upon his back, would result more satisfactorily than by the use
of any of the various appliances figured in the books.

\smallornament
\section*{Vienna and Its Hospital.}

\SectionStartWords{Our} readers will observe that our intelligent friends, \typo{Drs.}{Dr.'s} T.~C.\ and
L.~E.\ Miller, frequently refer to their experience and their professional
friends in Vienna, and it may be of interest to know that Vienna, the
capital of Austria, contains about 500,000 inhabitants.

The general Hospital was founded about one hundred years ago, by
the Emperor Joseph the II. It is the largest Hospital in the world,
the grounds embracing seven or eight squares, and is laid out in walks
shaded by rows of fine large trees.

Each square of ground is surrounded by a continuous line of buildings
three stories high, and in all they contain 3,000 beds, which, however,
are not always filled with patients, the usual number not averaging
more than from 2,300 to 2,500.

The clinics are held in this general Hospital, and usually there are
about 1,500 physicians and students in attendance on the clinics.

In the lying-in, or obstetrical department from 8,000 to 10,000 births
occur yearly.

In this city, with its University and Hospital, are presented admiral
and full opportunities for study and observation and we hope to present
to the readers of the \ThisJournal{College Journal} information in regard to
professional matters, as they occur in this great city, probably in each
number we issue.

\begin{figure}[H]
  \centering
  \includegraphics{pages/illustrations/bullet_arrow_divider.jpg}
\end{figure}

\chapter*[Book Notices.]{BOOK NOTICES.}

\footnotesize
\booktitle{The Effects of Climate on Tuberculous Diseases.}  By \ProperName{Edwin Lee}, M.~R.~C.~S., London. Being the Dissertation to which the Fiske Fund Prize was awarded June 6, 1855.

\booktitle{The Influence of Pregnancy on the Development of Tubercles.}  By \ProperName{Edward Warren}, \md, of Edenton, N.~C. Being the Dissertation to which the Fiske Fund Prize was awarded June 4, 1856. Philapelphia, \ProperName{Blanchard} and \ProperName{Lea}. 1857.
\plainbreak{1}

\normalsize

\SectionStartWords{Dr.\ Caleb Fiske}, of Rhode Island, bequeathed at his death to the
State Medical Society a fund of two thousand dollars, directing that
the annual income of that sum should be expended in premiums for
Medical Essays, on such subjects as should be designated by the\endinput
